% Options for packages loaded elsewhere
\PassOptionsToPackage{unicode}{hyperref}
\PassOptionsToPackage{hyphens}{url}
\documentclass[
]{article}
\usepackage{xcolor}
\usepackage[margin=1in]{geometry}
\usepackage{amsmath,amssymb}
\setcounter{secnumdepth}{-\maxdimen} % remove section numbering
\usepackage{iftex}
\ifPDFTeX
  \usepackage[T1]{fontenc}
  \usepackage[utf8]{inputenc}
  \usepackage{textcomp} % provide euro and other symbols
\else % if luatex or xetex
  \usepackage{unicode-math} % this also loads fontspec
  \defaultfontfeatures{Scale=MatchLowercase}
  \defaultfontfeatures[\rmfamily]{Ligatures=TeX,Scale=1}
\fi
\usepackage{lmodern}
\ifPDFTeX\else
  % xetex/luatex font selection
    \setmainfont[]{Times New Roman}
\fi
% Use upquote if available, for straight quotes in verbatim environments
\IfFileExists{upquote.sty}{\usepackage{upquote}}{}
\IfFileExists{microtype.sty}{% use microtype if available
  \usepackage[]{microtype}
  \UseMicrotypeSet[protrusion]{basicmath} % disable protrusion for tt fonts
}{}
\makeatletter
\@ifundefined{KOMAClassName}{% if non-KOMA class
  \IfFileExists{parskip.sty}{%
    \usepackage{parskip}
  }{% else
    \setlength{\parindent}{0pt}
    \setlength{\parskip}{6pt plus 2pt minus 1pt}}
}{% if KOMA class
  \KOMAoptions{parskip=half}}
\makeatother
\usepackage{graphicx}
\makeatletter
\newsavebox\pandoc@box
\newcommand*\pandocbounded[1]{% scales image to fit in text height/width
  \sbox\pandoc@box{#1}%
  \Gscale@div\@tempa{\textheight}{\dimexpr\ht\pandoc@box+\dp\pandoc@box\relax}%
  \Gscale@div\@tempb{\linewidth}{\wd\pandoc@box}%
  \ifdim\@tempb\p@<\@tempa\p@\let\@tempa\@tempb\fi% select the smaller of both
  \ifdim\@tempa\p@<\p@\scalebox{\@tempa}{\usebox\pandoc@box}%
  \else\usebox{\pandoc@box}%
  \fi%
}
% Set default figure placement to htbp
\def\fps@figure{htbp}
\makeatother
\setlength{\emergencystretch}{3em} % prevent overfull lines
\providecommand{\tightlist}{%
  \setlength{\itemsep}{0pt}\setlength{\parskip}{0pt}}
\usepackage{amsmath}
\usepackage{amssymb}
\usepackage{bbm}
\usepackage{bookmark}
\IfFileExists{xurl.sty}{\usepackage{xurl}}{} % add URL line breaks if available
\urlstyle{same}
\hypersetup{
  pdfauthor={Juanwu Lu},
  hidelinks,
  pdfcreator={LaTeX via pandoc}}

\title{\textbf{Fall 2024}\\
\textbf{Homework 4}\\
\textbf{STAT 656: Bayesian Data Analysis}}
\author{Juanwu Lu\footnote{College of Engineering, Purdue University,
  West Lafayette, IN, USA}}
\date{}

\begin{document}
\maketitle

\section{Gibbs Sampling for Probit
Model}\label{gibbs-sampling-for-probit-model}

Consider \(N\) pairs of observations \((x_{i},y_{i})\), where
\(x_{i}\in\mathbf{R}^{d}\) and \(y_{i}\in\left\{0, 1\right\}\). These
are generated as follows: \[
z_{i}\sim\mathcal{N}(x_{i}\beta,1)\quad y_{i}=\mathbbm{1}(z_{i}> 0),\quad i=1,\ldots,N 
\] This is thus the probit model fro the Lecture slides. Write \(X\),
\(Y\), and \(Z\) for the collection of \(x_{i}\), \(y_{i}\), and
\(z_{i}\)'s. We will place a \(\mathcal{N}(0,\Sigma_{b})\) prior on
\(\beta\in\mathbb{R}^{d}\), and draw samples from the posterior
\(p(\beta,Z\mid X,Y)\). Observe that as in regression settings, we are
really only interested in \(\beta\), the \(Z\)'s only auxiliary
variables to allow us to implement a Gibbs sampler. In other words, this
is an example of a \emph{data-augmentation algorithm}.

\begin{enumerate}
\def\labelenumi{\arabic{enumi}.}
\tightlist
\item
  (7 points) Derive and write down the conditional distribution
  \(p(Z\mid X,Y,\beta)\).
\item
  (7 points) Derive and write down the conditional distribution
  \(p(\beta\mid X,Y,Z)\)?
\item
  (1 points) Describe the overall Gibbs sampling algorithm, and how you
  would calculate the posterior mean and covariance of \(\beta\).
\item
  (20 points) Write an R function that implement the VB algorithm for
  the probit model given in the slides. This should accept as input a
  dataset \((X, Y)\) and return the approximations to the posterior
  distribution over \(\beta\) and the \(z\)'s. Recall how the algorithm
  proceeds: start with some arbitrary initial distribution \(q(\beta)\)
  over \(\beta\), use that to calculate the \(q(z_{i})\)'s, use those to
  calculate \(q(\beta)\), and repeat till these distributions stop
  changing.
\end{enumerate}

\textbf{Solution:}

\begin{enumerate}
\def\labelenumi{\arabic{enumi}.}
\item
  The conditional distribution \(Z\mid X, Y,\beta\) is derived as
  follows: \[
  \tag{1}
  \begin{aligned}
    \log{p}(Z|X,Y,\beta) &= \log{p}(\beta) + \sum\limits_{i=1}^{N}\log{p}(y_{i}\mid z_{i}) + \log{p}(z_{i}\mid x_{i}) + C_{1} \\
    &= -\frac{1}{2}\beta^{T}\Sigma_{b}^{-1}\beta + \sum\limits_{i=1}^{N}y_{i}\log\mathbbm{1}(z_{i}>0) + (1-y_{i})\log\mathbbm{1}(z_{i}\leq 0) - \frac{1}{2}\left(z_{i}-x_{i}\beta\right)^{2} + C_{2}
  \end{aligned}
  \] Therefore, the conditional distribution is given by: \[
  \tag{2}
  p(z_{i}|x_{i},y_{i},\beta)\propto\begin{cases}
    \mathbbm{1}(z_{i}>0)\mathcal{N}(x_{i}\beta,1) & \text{if } y_{i}=1, \\
    \mathbbm{1}(z_{i}\leq 0)\mathcal{N}(x_{i}\beta,1) & \text{otherwise.}
  \end{cases}
  \]
\item
  Based on the results from equation (1), the conditional distribution
  of \(\beta\mid X,Y,Z\) is derived by: \[
  \begin{aligned}
  \log{p}(\beta\mid X,Y,Z) &= -\frac{1}{2}\sum\limits_{i=1}^{d}z_{i}^{2}-2z_{i}(x_{i}\beta) + x_{i}\beta\beta^{\intercal}x_{i}^{\intercal} + \beta^{\intercal}\Sigma_{b}^{-1}\beta + C_{3} \\
  &= -\frac{1}{2}\beta^{\intercal}\left(\sum\limits_{i=1}^{N}x_{i}^{\intercal}x_{i}+\Sigma_{b}^{-1}\right)\beta - 2\sum\limits_{i=1}^{N}z_{i}x_{i}\beta + C_{4} \\
  &= -\frac{1}{2}\left(\beta-\Sigma_{\beta\mid X,Y,Z}\sum\limits_{i=1}^{N}z_{i}x_{i}^{\intercal}\right)^{\intercal}\Sigma_{\beta\mid X,Y,Z}^{-1}\left(\beta-\Sigma_{\beta\mid X,Y,Z}\sum\limits_{i=1}^{N}z_{i}x_{i}^{\intercal}\right) + C_{5}, \\
  \text{where}\quad & \Sigma_{\beta\mid X,Y,Z} = \left(\sum\limits_{i=1}^{N}x_{i}^{\intercal}x_{i}+\Sigma_{b}^{-1}\right)^{-1}
  \end{aligned}
  \] Therefore, the conditional distribution is a multivariate normal
  distribution \[
  \tag{2}
  p(\beta\mid X,Y,Z)\sim\mathcal{N}\left(\left(\sum\limits_{i=1}^{N}x_{i}^{\intercal}x_{i}+\Sigma_{b}^{-1}\right)^{-1}\sum\limits_{i=1}^{N}z_{i}x_{i}^{\intercal}, \left(\sum\limits_{i=1}^{N}x_{i}^{\intercal}x_{i}+\Sigma_{b}^{-1}\right)^{-1}\right)
  \]
\item
  The Gibbs sampling algorithm is as follows:
\end{enumerate}

\begin{itemize}
\tightlist
\item
  Initialize \(\beta\) and \(Z\).
\item
  Repeat the following steps until convergence or a maximum number of
  iterations:

  \begin{itemize}
  \tightlist
  \item
    Sample \(Z\mid X,Y,\beta\) from the conditional distribution
    \(p(Z\mid X,Y,\beta)\).
  \item
    Sample \(\beta\mid X,Y,Z\) from the conditional distribution
    \(p(\beta\mid X,Y,Z)\). Therefore, with every sampled collection of
    \(Z\), we can calculate the posterior mean and covariance of
    \(\beta\) according to the results equation (\(2\)):
  \end{itemize}
\item
  Calculate the covariance by
  \(\left(\sum\limits_{i=1}^{d}x_{i}^{\intercal}x_{i}+\Sigma_{b}^{-1}\right)^{-1}\).
\item
  Calculate the mean by
  \(\left(\sum\limits_{i=1}^{d}x_{i}^{\intercal}x_{i}+\Sigma_{b}^{-1}\right)^{-1}\sum\limits_{i=1}^{d}z_{i}x_{i}^{\intercal}\),
  where \(z_{i}\) are samples from the first step of Gibbs sampling
  iterations.
\end{itemize}

\begin{enumerate}
\def\labelenumi{\arabic{enumi}.}
\setcounter{enumi}{3}
\tightlist
\item
  The mean-field variational family of the original model is given by
  \(q(\beta,z)=q(\beta)\prod\limits_{i=1}^{\)
\end{enumerate}

\pagebreak

\section{Effect of Computer-generated
Reminders}\label{effect-of-computer-generated-reminders}

\end{document}
